\begin{figure}[t]
\centering
\resizebox{\textwidth}{!}{
\begin{tikzpicture}[
  layer/.style={draw, rounded corners=1.8pt, minimum height=18pt,
                minimum width=26pt, font=\footnotesize, inner sep=2.5pt},
  preloop/.style={layer, fill=gray!10, draw=gray!55!black, line width=0.45pt},
  loopwin/.style={layer, fill=blue!12, draw=blue!55!black, line width=0.8pt},
  postloop/.style={layer, fill=gray!10, draw=gray!55!black, line width=0.45pt},
  arrow/.style={-{Stealth[length=3pt]}, line width=0.6pt,
                shorten >=1pt, shorten <=1pt},
  loopback/.style={-{Stealth[length=4pt]}, line width=0.95pt,
                   draw=blue!55!black},
  perlayerloop/.style={-{Stealth[length=3pt]}, line width=0.8pt,
                       draw=blue!55!black},
  font=\footnotesize,
  node distance=4pt,
]

\begin{scope}[shift={(0,1.6)}]

\node[preloop] (L0a)   at (0.0, 0)  {$L_0$};
\node[preloop, right=6pt of L0a] (L1a)  {$L_1$};
\node[right=6pt of L1a]      (dotsAa){$\cdots$};
\node[preloop, right=6pt of dotsAa] (Lam1a) {$L_{a{-}1}$};

\node[loopwin, right=11pt of Lam1a]  (Laa)   {$L_a$};
\node[right=6pt of Laa]              (dotsBa){$\cdots$};
\node[loopwin, right=6pt of dotsBa]  (Lba)   {$L_b$};

\node[postloop, right=11pt of Lba]   (Lbp1a) {$L_{b{+}1}$};
\node[right=6pt of Lbp1a]            (dotsCa){$\cdots$};
\node[postloop, right=6pt of dotsCa] (LNa)   {$L_{N{-}1}$};

\node[font=\footnotesize\itshape, anchor=east]
   at ($(L0a.west) + (-0.85, 0)$) {(a) block-mode};

\draw[arrow] ($(L0a.west) + (-0.32,0)$) node[left]{$x$} -- (L0a);
\draw[arrow] (L0a) -- (L1a); \draw[arrow] (L1a) -- (dotsAa);
\draw[arrow] (dotsAa) -- (Lam1a); \draw[arrow] (Lam1a) -- (Laa);
\draw[arrow] (Laa) -- (dotsBa); \draw[arrow] (dotsBa) -- (Lba);
\draw[arrow] (Lba) -- (Lbp1a); \draw[arrow] (Lbp1a) -- (dotsCa);
\draw[arrow] (dotsCa) -- (LNa);
\draw[arrow] (LNa) -- ($(LNa.east) + (0.32,0)$) node[right]{$\hat f(x)$};

\draw[loopback]
   ([yshift=2pt]Lba.north) -- ++(0,0.32) -| ([yshift=2pt]Laa.north)
   node[pos=0.25, above=-1pt, font=\footnotesize\bfseries,
        color=blue!55!black] {$\times K$};
\end{scope}

\begin{scope}[shift={(0,0)}]

\node[preloop] (L0b)   at (0.0, 0)  {$L_0$};
\node[preloop, right=6pt of L0b] (L1b)  {$L_1$};
\node[right=6pt of L1b]      (dotsAb){$\cdots$};
\node[preloop, right=6pt of dotsAb] (Lam1b) {$L_{a{-}1}$};

\node[loopwin, right=11pt of Lam1b]  (Lab)   {$L_a$};
\node[right=6pt of Lab]              (dotsBb){$\cdots$};
\node[loopwin, right=6pt of dotsBb]  (Lbb)   {$L_b$};

\node[postloop, right=11pt of Lbb]   (Lbp1b) {$L_{b{+}1}$};
\node[right=6pt of Lbp1b]            (dotsCb){$\cdots$};
\node[postloop, right=6pt of dotsCb] (LNb)   {$L_{N{-}1}$};

\node[font=\footnotesize\itshape, anchor=east]
   at ($(L0b.west) + (-0.85, 0)$) {(b) layer-mode};

\draw[arrow] ($(L0b.west) + (-0.32,0)$) node[left]{$x$} -- (L0b);
\draw[arrow] (L0b) -- (L1b); \draw[arrow] (L1b) -- (dotsAb);
\draw[arrow] (dotsAb) -- (Lam1b); \draw[arrow] (Lam1b) -- (Lab);
\draw[arrow] (Lab) -- (dotsBb); \draw[arrow] (dotsBb) -- (Lbb);
\draw[arrow] (Lbb) -- (Lbp1b); \draw[arrow] (Lbp1b) -- (dotsCb);
\draw[arrow] (dotsCb) -- (LNb);
\draw[arrow] (LNb) -- ($(LNb.east) + (0.32,0)$) node[right]{$\hat f(x)$};

\draw[perlayerloop]
   ([xshift=-3pt, yshift=2pt]Lab.north)
   .. controls +(-0.10, 0.26) and +(0.10, 0.26) ..
   ([xshift=3pt, yshift=2pt]Lab.north);
\node[font=\scriptsize\bfseries, color=blue!55!black, anchor=south]
   at ($(Lab.north) + (0,0.21)$) {$\times K$};

\draw[perlayerloop]
   ([xshift=-3pt, yshift=2pt]Lbb.north)
   .. controls +(-0.10, 0.26) and +(0.10, 0.26) ..
   ([xshift=3pt, yshift=2pt]Lbb.north);
\node[font=\scriptsize\bfseries, color=blue!55!black, anchor=south]
   at ($(Lbb.north) + (0,0.21)$) {$\times K$};

\node[font=\scriptsize\bfseries, color=blue!55!black, anchor=south]
   at ($(dotsBb.north) + (0,0.24)$) {$\times K$};

\draw[decorate, decoration={brace,amplitude=3pt,mirror,raise=2pt},
      semithick, color=blue!55!black]
   (Lab.south west) -- (Lbb.south east)
   node[midway, below=4pt, font=\footnotesize\bfseries,
        color=blue!55!black]
   {loop window $g$};

\draw[decorate, decoration={brace,amplitude=3pt,mirror,raise=2pt},
      thin, color=gray!55!black]
   (L0b.south west) -- (Lam1b.south east)
   node[midway, below=4pt, font=\footnotesize, color=gray!50!black]
   {pre-loop};

\draw[decorate, decoration={brace,amplitude=3pt,mirror,raise=2pt},
      thin, color=gray!55!black]
   (Lbp1b.south west) -- (LNb.south east)
   node[midway, below=4pt, font=\footnotesize, color=gray!50!black]
   {post-loop};

\end{scope}
\end{tikzpicture}
}
\vspace{-2pt}
\caption{\textbf{Training-free looped transformer wrapper, two iteration modes.} A frozen checkpoint is augmented at inference by re-applying a contiguous mid-block $g = L_b\circ\cdots\circ L_a$ for $K$ iterations before resuming the post-loop layers. No weights are changed and no new parameters are introduced; the cost is $(b-a+1)(K-1)$ extra forward passes through the loop window. \textbf{(a) Block-mode} iterates the whole window $K$ times, $(L_b\circ\cdots\circ L_a)^K$. \textbf{(b) Layer-mode} iterates each window layer $K$ times before passing on, $L_b^K\circ\cdots\circ L_a^K$, which is the safer default on Mixture-of-Experts backbones because it pins per-layer expert routing across iterations (Section~\ref{sec:layer-mode}). Section~\ref{sec:strategy} interprets $g^{(K)}$ as sub-stepping the residual ODE the network already integrates with single-step forward Euler.}
\label{fig:method}
\end{figure}


\section{Training-free looped transformers}
\label{sec:method}

\subsection{Loop wrapper}
\label{sec:wrapper}

% \begin{algorithm}[t]
% \caption{Explicit \HiB{\textcolor{blockdark}{\textbf{block-mode}}} \HiL{\textcolor{layerdark}{\textbf{layer-mode}}} Runge--Kutta forward pass.}
% \label{alg:rk-unified}
% \begin{algorithmic}[1]
% \Require $x$, range $[a,b]$, number of stages $s$, Runge--Kutta coefficients $\{a_{ij}\}_{1\le j<i\le s},\{b_i\}_{i=1}^{s}$
% \State $x \gets \textsc{PreLoop}(x)$ \Comment{embedding, pre-block, etc.}
% \State \HiB{$y_0 \gets x$} \Comment{\HiB{initial state for the block}}
% \State \HiL{\textbf{for} $\ell = a, \ldots, b$ \textbf{do}} \Comment{\HiL{outer loop over each layer}}
% \State \HiL{\quad $y_0 \gets x$} \Comment{\HiL{initial state for every layer}}
% \State \quad \textbf{for} $i = 1, \ldots, s$ \textbf{do} \Comment{iterate over the $s$ Runge--Kutta stages}
% \State \quad\quad $y_i \gets y_0 + \sum_{j<i} a_{ij}k_j$ \Comment{add weighted increments}
% \State \quad\quad $k_i \gets$ \HiB{$(L_b \circ \cdots \circ L_a)(y_i) - y_i$}\,\HiL{$L_\ell(y_i) - y_i$} \Comment{residual of the \HiB{full block}\,\HiL{single layer}}
% \State \quad \textbf{end for}
% \State \quad $x \gets y_0 + \sum_{i=1}^{s} b_ik_i$ \Comment{$s^\text{th}$ order Runge--Kutta update}
% \State \HiL{\textbf{end for}}
% \State \Return $\textsc{PostLoop}(x)$ \Comment{layers after the loop block, output head}
% \end{algorithmic}
% \end{algorithm}

\begin{algorithm}[t]
\caption{\HiB{\textbf{\textcolor{blockdark}{Block-mode}}} Runge--Kutta forward pass.}
\label{alg:rk-block}
\begin{algorithmic}[1]
\Require $x$, block range $[a,b]$, number of stages $s$, Runge--Kutta coefficients $\{a_{ij}\}_{1\leq j<i\leq s},\{b_i\}_{i=1}^s$
\State $x \gets \textsc{PreLoop}(x)$ \Comment{embedding, pre-block, etc.}
\State \HiB{$y_0 \gets x$} \Comment{\HiB{initial state for the block}}
\For{$i = 1,\dots,s$} \Comment{iterate over the $s$ Runge--Kutta stages}
  \State $y_i \gets y_0 + \sum_{j<i} a_{ij}k_j$ \Comment{add weighted increments}
  \State \HiB{$k_i \gets (L_b\circ\cdots\circ L_a)(y_i) - y_i$} \Comment{\HiB{residual of the full composed block}}
\EndFor
\State {$x \gets y_0 + \sum_{i=1}^{s} b_ik_i$} \Comment{$s^\text{th}$ order Runge--Kutta update}
\State \Return $\textsc{PostLoop}(x)$ \Comment{layers after the loop block, output head}
\end{algorithmic}

\end{algorithm}
\begin{algorithm}[t]
\caption{\HiL{\textbf{\textcolor{layerdark}{Layer-mode}}} Runge--Kutta forward pass.}
\label{alg:rk-layer}
\begin{algorithmic}[1]
\Require $x$, block range $[a,b]$, number of stages $s$, Runge--Kutta coefficients $\{a_{ij}\}_{1\leq j<i\leq s},\{b_i\}_{i=1}^s$
\State $x \gets \textsc{PreLoop}(x)$ \Comment{embedding, pre-block, etc.}
\State \HiL{\textbf{for} $\ell = a, \ldots, b$ \textbf{do}} \Comment{\HiL{outer loop over each layer}}
  \State\quad \HiL{$y_0 \gets x$} \Comment{\HiL{initial state for every layer}}
  \State \quad \textbf{for} $i = 1, \ldots, s$ \textbf{do} \Comment{iterate over the $s$ Runge--Kutta stages}
    \State \quad \quad $y_i \gets y_0 + \sum_{j<i} a_{ij}k_j$ \Comment{add weighted increments}
    \State \quad \quad\HiL{$k_i \gets L_\ell(y_i) - y_i$} \Comment{\HiL{residual of a single layer $L_\ell$}}
\State \quad \textbf{end for}
  \State \quad {$x \gets y_0 + \sum_{i=1}^{s} b_ik_i$} \Comment{$s^\text{th}$ order Runge--Kutta update}
\State \HiL{\textbf{end for}}
\State \Return $\textsc{PostLoop}(x)$ \Comment{layers after the loop block, output head}
\end{algorithmic}
\end{algorithm}

Let $f = L_{N-1} \circ \cdots \circ L_0$ denote a pretrained transformer with $N$ decoder layers \citep{vaswani2017attention,qwen3,llama3,deepseekv2,moonlight}, where $L_i$ maps a residual stream of shape $\mathbb{R}^{T \times d}$ to itself. Choose a contiguous \emph{loop window} indexed by $[a, b]$ with $0 \le a \le b \le N-1$ and a loop count $K \ge 1$. The window induces an operator \begin{equation}
g := L_b \circ L_{b-1} \circ \cdots \circ L_a,
\qquad g : \mathbb{R}^{T \times d} \to \mathbb{R}^{T \times d}.
\label{eq:block-op}
\end{equation}
The looped wrapper splits the network into pre-loop layers $0, \ldots, a-1$, the looped middle, and post-loop layers $b+1, \ldots, N-1$:
\begin{equation}
\hat f(x) =
\underbrace{\bigl(L_{N-1} \circ \cdots \circ L_{b+1}\bigr)}_{\text{post-loop}}
\;\circ\; g^{(K)} \;\circ\;
\underbrace{\bigl(L_{a-1} \circ \cdots \circ L_0\bigr)}_{\text{pre-loop}}(x),
\label{eq:wrapper}
\end{equation}
where $g^{(K)} : \mathbb{R}^{T \times d} \to \mathbb{R}^{T \times d}$ is a $K$-step iteration of $g$. When $a = 0$ or $b = N-1$, the corresponding boundary composition is the identity. The map $g^{(K)}$ depends on an \emph{iteration mode} (Section~\ref{sec:layer-mode}) and a \emph{loop strategy} (Section~\ref{sec:strategy}).

\subsection{Block-mode versus layer-mode iteration}
\label{sec:layer-mode}

Given the window operator $g = L_b \circ \cdots \circ L_a$ in \eqref{eq:block-op}, there are two natural realizations of $g^{(K)}$, \begin{align}
\text{block-mode:}\quad
   &g^{(K)}_{\operatorname{blk}}(x) = (L_b \circ \cdots \circ L_a)^{K}(x), \label{eq:block-mode}\\
\text{layer-mode:}\quad
   &g^{(K)}_{\operatorname{lyr}}(x) = L_b^{K} \circ L_{b-1}^{K} \circ \cdots \circ L_a^{K}(x).
     \label{eq:layer-mode}
\end{align}
In other words, block-mode iterates the entire window as one unit, whereas layer-mode iterates each layer before passing to the next (Figure~\ref{fig:method}). For dense models \citep{bae2024relaxed,bae2025mor}, the two yield broadly similar quality. For mixture-of-experts (MoE) layers \citep{deepseekv2,moonlight,csordas2024moeut}, however, block-mode becomes unstable, since at iteration $i$ the gating network of every MoE layer in the window sees a slightly perturbed hidden state and routes to a different subset of experts than at iteration $i-1$, and the accumulated routing-induced noise eventually overpowers the intended refinement. On the other hand, layer-mode iteration computes the gating decision once and applies the same expert mixture $K$ times \citep{csordas2024moeut}, avoiding this failure mode and making it the correct default for MoE backbones.

\subsection{Loop strategies}
\label{sec:strategy}

We first motivate several loop strategies with a structural observation. A standard pre-norm transformer layer $L$ \citep{vaswani2017attention} implements
\begin{equation}
L(x) = x + \operatorname{Attn}(\operatorname{LN}_1(x))+\operatorname{MLP}(\operatorname{LN}_2(x + \operatorname{Attn}(\operatorname{LN}_1(x)))).
\label{eq:block}
\end{equation}

The right-hand side is the layer's input plus a residual update. Generalizing this directly to the loop window operator $g$, which may be a single layer $g = L_i$ (the layer-mode case) or a contiguous composition $g = L_b \circ \cdots \circ L_a$ (the block-mode case), we define the \emph{window residual field}
\begin{equation}
F_g(x) := g(x) - x.
\label{eq:F}
\end{equation}

By construction, $g(x) = x + F_g(x)$, which is exactly a forward Euler step with step size $h{=}1$
\begin{equation}
    x_1 = x_0 + h \cdot F_g(x_0)
\end{equation}
on the autonomous ODE \citep{bai2019deep,bai2020multiscale}
\begin{equation}
\dot{x}=F_g(x). 
\label{eq:ode}
\end{equation}
This holds regardless of whether the window contains one layer or many, and the two modes differ only in what $F_g$ unfolds to. In layer-mode, $F_g$ is just the single-layer residual field \[ F_g(x) = \operatorname{Attn}_i(\operatorname{LN}^i_1(x)) + \operatorname{MLP}_i(\operatorname{LN}^i_2(\cdot)). \]
In block-mode, telescoping the unrolled chain (proof in Appendix~\ref{app:block-unroll}) gives
\begin{equation}
F_g(x) = \sum_{i=a}^{b} F_{L_i}(y_i(x)),
\qquad y_a(x) := x,\quad y_{i+1}(x) := y_i(x) + F_{L_i}(y_i(x)),
\label{eq:block-effective-field}
\end{equation}
where $F_{L_i}(z) := L_i(z) - z$ is the layer residual field of layer $i$. In both modes, the post-loop layers $L_{b+1}, \ldots, L_{N-1}$ are trained to receive the approximation of $x(t{=}1)$ on $F_g$ implicitly produced by one application of $g$ and not the trajectory at any other time.

\begin{figure}[t]
\centering
\includegraphics[width=\linewidth]{figures/toy_mlp_contour_2.pdf}
\caption{\textbf{RK integration vs.\ naive looping.} A tiny MLP $\operatorname{pre}(\mathbb{R}^4\to\mathbb{R}^2)\to$ block of 3 residual layers $(\mathbb{R}^2\to\mathbb{R}^2) \to \operatorname{post}(\mathbb{R}^2\to\mathbb{R}^2)$ is trained end-to-end on a 2-D regression target. Each panel fixes $K \in \{2,4,8\}$ and shows, over a $220{\times}220$ grid in the post-block hidden state $\mathbf{z}$, the median test loss $L(\mathbf{z}) = \operatorname{med}_i \norm{\operatorname{post}(\mathbf{z}) - y_i}^2$ (log scale, colored) together with the test-set scatters obtained by naive looping (red circles) and $K$-stage RK integration (purple squares). The trained baseline induces a central low-loss valley (blue); our endpoints stay clustered in that valley at every $K$, while naive $K$-loop endpoints drift increasingly outward into high-loss regions (yellow) as $K$ increases.}
\label{fig:toy-mlp-contour}
\end{figure}


Naively looping $g$ for $K$ rounds, i.e. $x \leftarrow g(x)$ applied $K$ times \citep{dehghani2019universal,lan2020albert,liu2024looped,yang2025looped,saunshi2025reasoning,xu2025expressive,geiping2025scaling,zhu2025ouro,vonoswald2023transformers,ahn2024transformers,mahankali2024one,gatmiry2024looped,chen2024bypassing}, is a $K$-step forward Euler integration of the window residual field with step size $h{=}1$, which approximates $x(t{=}K)$. But the post-loop layers are not trained to receive the trajectory at $t{=}K$, and empirically, naive $K$-fold looping degrades performance almost universally (Section~\ref{sec:experiments}). Figure~\ref{fig:toy-mlp-contour} illustrates the effect of naive looping vs.\ sub-stepping on a tiny end-to-end trainable network with a $2$-D bottleneck, where the input to the post-loop layers can be directly plotted. The damped sub-step endpoints stay clustered in the trained low-loss valley, while the naive $K$-loop endpoints drift into high-loss regions. Numerically, when averaged over the test set, naive $K{=}2$ looping increases the MSE to 2.88 versus the baseline and sub-step values of 0.015 and 0.36, respectively ($\approx 8\times$ gap); at $K{=}8$ this degrades to a MSE of 335 versus the sub-step value of 1.04 ($\approx 320\times$ gap).

The principled goal of $g^{(K)}$ is therefore not to advance integration to $t{=}K$, but to \emph{better approximate the same endpoint} $x(t{=}1)$ that the unmodified network already targets. To address this challenge, classical numerical analysis suggests to sub-step the same total integration time $[0,1]$ at finer resolution $h = 1/K$. Performing $K$ Euler steps of size $h{=}1/K$ on \eqref{eq:ode} gives the damped update
\begin{equation}
x_{k+1}=\left(1-\frac{1}{K}\right)x_k+\frac{1}{K}g(x_k),
\label{eq:sub-step-euler}
\end{equation}
which converges to the true $x(t{=}1)$ at order $O(1/K)$ and strictly improves upon the single-shot $h{=}1$ Euler step that the network implicitly implements at every layer. Because equations~\eqref{eq:F}--\eqref{eq:ode} are mode-agnostic, this principle applies in both layer-mode (sub-stepping the per-layer field $F_g = F_{L_i}$) and in block-mode (sub-stepping the composite field $F_g$ in~\eqref{eq:block-effective-field}).

More generally, we can consider numerical integration strategies beyond the forward Euler method to approximate $x(t{=}1)$. In particular, the damped Euler update~\eqref{eq:sub-step-euler} can be replaced by the $s$-stage explicit Runge--Kutta integrator
\begin{equation}
x_1 = x_0 + h\sum_{i=1}^{s} b_i k_i,
\label{eq:rk-update}
\end{equation}
where the $s$ stages are given by
\begin{equation}
k_1 = F_g(x_0),\qquad k_2 = F_g(x_0 + h (a_{21} k_1)),\qquad\dots,\qquad k_s = F_g\Par{x_0 + h\sum_{j=1}^{s-1} a_{sj}k_j},
\label{eq:rk-stages}
\end{equation}
with the coefficients $\{a_{ij}, b_i\}$ specified by a Butcher tableau. Algorithms~\ref{alg:rk-block} and~\ref{alg:rk-layer} summarize this family of methods for block-mode and layer-mode implementations, respectively, where the step size is chosen to be $h{=}1$.

\begin{algorithm}[t]
\caption{Explicit \HiB{\textcolor{blockdark}{\textbf{block-mode}}} \HiL{\textcolor{layerdark}{\textbf{layer-mode}}} Runge--Kutta with Butcher tableau \eqref{eq:rk-euler-chain-aij}–\eqref{eq:rk-beta-bi}.}
\label{alg:rk_beta}
\begin{algorithmic}[1]
\Require $x$, range $[a,b]$, number of stages $K$, anchor weight $\beta\in[0,1]$
\State $x \gets \textsc{PreLoop}(x)$ \Comment{embedding, pre-block, etc.}
\State \HiB{$\tilde{x} \gets (L_b \circ \cdots \circ L_a)(x)$} \Comment{\HiB{anchor (block mode)}}
\State \HiL{\textbf{for} $\ell = a, \ldots, b$ \textbf{do}} \Comment{\HiL{outer loop over each layer}}
\State \quad \HiL{$\tilde{x} \gets L_i(x)$} \Comment{\HiL{anchor (layer mode)}}
\State \quad \textbf{for} $k = 1, \ldots, K-1$ \textbf{do} \Comment{iterate over the $K$ Runge--Kutta stages}
\State \quad\quad \HiB{$y \gets (L_b \circ \cdots \circ L_a)(x)$} \HiL{$y \gets L_\ell(x)$}
\State \quad\quad $x \gets x + \frac{1}{K}(y - x)$ \Comment{Runge--Kutta update}
\State \quad \textbf{end for}
\State \quad $x \gets \beta\tilde{x} + (1-\beta)x$ \Comment{anchor combination}
\State \HiL{\textbf{end for}}
\State \Return $\textsc{PostLoop}(x)$
\end{algorithmic}
\end{algorithm}

As a particularly effective example, we choose a specific Butcher tableau so that the RK output becomes an interpolation between the base output and the $K$-step damped Euler updates following~\eqref{eq:sub-step-euler}. Concretely, we set
\begin{equation}
h=1,\quad s=K,\quad\text{and}\quad a_{ij}=\frac{1}{K}\ind(j<i),
\label{eq:rk-euler-chain-aij}
\end{equation}
so that the $i^\text{th}$ stage evaluates $F_g$ at the point
\begin{equation}\label{eq:yi}
y_i=x_0+\frac{1}{K}\sum_{j=1}^{i-1}k_j
    = F^{i-1}(x_0),\qquad k_i=F_g(y_i),
\end{equation}
where $F(x)\defeq x+\frac{1}{K}F_g(x)$ is the damped Euler update~\eqref{eq:sub-step-euler}. We then choose $\beta\in[0,1]$ and set the output weights as
\begin{equation}
b_1=\beta+\frac{1-\beta}{K},
\qquad
b_i=\frac{1-\beta}{K},\quad i=2,\ldots,K,
\label{eq:rk-beta-bi}
\end{equation}
which are nonnegative and sum to one. Under these coefficients, the RK output satisfies the identity
\begin{equation}
x_1=x_0+\sum_{i=1}^{K} b_i k_i=\beta g(x_0)+(1-\beta)F^{K}(x_0),
\label{eq:rk-beta-equivalence}
\end{equation}
which leads to an efficient implementation (Algorithm~\ref{alg:rk_beta}). The proof is provided in Appendix~\ref{app:block-unroll}.

Algorithm~\ref{alg:rk_beta} can be interpreted as a RK method with a front-loaded quadrature rule controlled by $\beta$. In the extreme cases, $\beta=0$ recovers the $K$-substep forward Euler endpoint, while $\beta=1$ preserves the original output $g(x_0)$. For intermediate values of $\beta$, the method effectively biases the trajectory toward the trained one-step endpoint by placing greater weight on the initial residual direction.

\begin{table}[t]
\centering
\caption{\textbf{Iteration strategies for numerical integration of the ODE \eqref{eq:ode}}.}
\small
\setlength{\tabcolsep}{8pt}
\renewcommand{\arraystretch}{1.5}
\resizebox{\textwidth}{!}{%
\begin{tabular}{@{}l l c@{}}
\toprule
\textbf{Strategy} & \textbf{Update Rule} & \textbf{Forward Passes} \\
\midrule
\multicolumn{3}{@{}l}{\textit{Naive iteration \& forward Euler}}\\
\cmidrule(lr){1-3}
Naive Loop
  & $x_{k+1} = g(x_k)$
  & $K$ \\
Forward Euler~\citep{bai2019deep,bai2020multiscale}
  & $x_{k+1} = x_k + \frac{1}{K}F_g(x_k)$
  & $K$ \\
\midrule
\multicolumn{3}{@{}l}{\textit{Higher-order Runge--Kutta}~\citep{butcher2008numerical}}\\
\cmidrule(lr){1-3}
Midpoint (RK2)
  & $x_{k+1} = x_k + hF_g(x_k + \frac{h}{2}F_g(x_k))$
  & $2K$ \\
Heun (RK2)
  & $x_{k+1} = x_k + \frac{h}{2}(F_g(x_k) + F_g(x_k + hF_g(x_k)))$
  & $2K$ \\
RK4
  & $x_{k+1} = x_k + \frac{h}{6}(\kappa_1 + 2\kappa_2 + 2\kappa_3 + \kappa_4)$
  & $4K$ \\
\midrule
\multicolumn{3}{@{}l}{\textit{Fixed-point accelerators}}\\
\cmidrule(lr){1-3}
Heavy-ball ($\alpha,\beta$)~\citep{polyak1964heavyball}
  & $x_{k+1} = x_k + \alpha F_g(x_k) + \beta(x_k - x_{k-1})$
  & $K$ \\
Anderson ($m,\beta$)~\citep{walker2011anderson}
  & $x_{k+1} = (1-\beta)(x_k - (\Delta X_k)\gamma_k^{\star}) + \beta(g(x_k) - (\Delta F_k)\gamma_k^{\star})$
  & $K$ \\
Aitken $\Delta^2$~\citep{aitken1926}
  & $x_{k+1} = x_k - \frac{(d_{1,k})^2}{d_{2,k}}$ \;\;(per-coordinate)
  & $K$ \\
Uniform Loop~\citep{lys2026inner}
  & $x_{k+1} = g\!\left(\frac{1}{k+1}\sum_{i=0}^{k} x_i\right)$
  & $K$ \\
\bottomrule
\end{tabular}}
\scriptsize\raggedright
\emph{Notation.} RK4: $\kappa_1 := F_g(x_k)$,
$\kappa_2 := F_g(x_k + \frac{h}{2}\kappa_1)$,
$\kappa_3 := F_g(x_k + \frac{h}{2}\kappa_2)$,
$\kappa_4 := F_g(x_k + h\kappa_3)$.\\
Anderson:
$\gamma_k^{\star} := \argmin_\gamma\norm{f_k -(\Delta F_k)\gamma}_2$ with $f_i := g(x_i) - x_i$ and $\Delta X_k, \Delta F_k$ the matrices of the last $m$ residual / state increments, respectively \citep{walker2011anderson}.\\
Aitken: $d_{1,k} := g(x_k) - x_k$, $d_{2,k} := g(g(x_k)) - 2g(x_k) + x_k$, applied per coordinate; the Steffensen safeguard clips $|x_{k+1} - x_k| \le |d_{1,k}|$ and requires $K$ even.
\label{tab:strategies}
\vspace{-1em}
\end{table}

Beyond Runge--Kutta methods, there exists a wide range of numerical integration schemes that have been extensively studied in the ODE literature \citep{walker2011anderson,bai2019deep}. These methods are often designed to address specific properties of the underlying dynamics, and we list several prominent examples in Table~\ref{tab:strategies}. These algorithms, which differ in order, accelerator, and step size choices, will form the basis for seven strategies that we evaluate.
